Questo lavoro presenta un modello d'ordine ridotto per l'interpolazione parametrica di flussi reattivi basato sulla Higher-Order Singular Value Decomposition (HOSVD) combinata con la Gaussian Process Regression (GPR). Rispetto al classico approccio POD+GPR, il metodo proposto sfrutta la struttura tensoriale dei dati di addestramento organizzati su una griglia cartesiana nello spazio dei parametri, disaccoppiando i modi parametrici da quelli spaziali. Due GPR univariati sostituiscono il singolo GPR bivariato, riducendo la dimensione effettiva del problema di regressione. Gli iperparametri dei kernel sono selezionati tramite una ricerca a griglia su un insieme di validazione dedicato. Il metodo è validato su un dataset di 100 simulazioni CFD di una fiamma a diffusione di idrogeno turbolenta, parametrizzate dal numero di Reynolds e dalla frazione di massa del combustibile su una griglia $10\times10$. HOSVD+GPR supera sistematicamente POD+GPR in termini di errore $\ell^2$ relativo per tutti e sei i campi esaminati (temperatura, CH$_4$, O$_2$, CO, CO$_2$, H$_2$O), con vantaggi più marcati nei punti di test in cui entrambi i parametri non sono presenti nella griglia di addestramento.
